% Pattern: proof/attach/label. A theorem of the introduction is restated with
% the same label before its proof, and the head of the proof refers to the
% label: the proof proves the restatement. A second head says that a lemma
% implies the theorem: the proof proves the theorem.
% Seen in the gold pilot (set E).
% want theorem-1 statement/theorem 1 status=none labels=main
% want lemma-1 statement/lemma 2 status=none labels=aux
% want theorem-2 statement/theorem 3 status=here labels=main
% want proof-1 proof - proves=theorem-2
% want theorem-3 statement/theorem 4 status=here labels=other
% want proof-2 proof - proves=theorem-3
\documentclass{article}
\newtheorem{theorem}{Theorem}
\newtheorem{lemma}[theorem]{Lemma}
\begin{document}
\begin{theorem}\label{main}
A statement.
\end{theorem}
\begin{lemma}\label{aux}
A helper.
\end{lemma}
Some text between.

\begin{theorem}\label{main}
A statement.
\end{theorem}
\begin{proof}[Proof of Theorem~\ref{main}]
A step.
\end{proof}
\begin{theorem}\label{other}
Another statement.
\end{theorem}
Some more text.

\begin{proof}[Proof that Lemma~\ref{aux} implies Theorem~\ref{other}]
Another step.
\end{proof}
\end{document}
